% Part: first-order-logic
% Chapter: axiomatic-deduction
% Section: proof-theoretic-notions

\documentclass[../../../include/open-logic-section]{subfiles}

\begin{document}

\iftag{FOL}
      {\olfileid[id]{fol}{axd}{ptn}}
      {\olfileid[id]{pl}{axd}{ptn}}

\olsection{Gagasan Teori-Bukti}

\begin{explain}
Sebagaimana kita telah mendefinisikan sejumlah gagasan semantik penting
(\iftag{FOL}{validitas}{tautologi}, perikutan, keterpenuhan), kini kita
mendefinisikan \emph{gagasan teori-bukti} yang bersesuaian.
Gagasan-gagasan ini tidak didefinisikan dengan mengacu pada terpenuhinya
!!{sentence}s di dalam \iftag{FOL}{!!{structure}s}{valuasi},
melainkan dengan mengacu pada !!{derivability} atau
!!{nonderivability} formula-formula tertentu. Penemuan bahwa
gagasan-gagasan tersebut berimpit merupakan penemuan penting.
Keberimpitan itu adalah isi \emph{teorema kesahihan} dan
\emph{teorema kelengkapan}.
\end{explain}

\begin{defn}[!!^{derivability}]
!!^a{formula}~$!A$
\emph{!!{derivable}} dari $\Gamma$, ditulis $\Gamma \Proves !A$, jika
terdapat !!a{derivation} dari~$\Gamma$ yang berakhir dengan~$!A$.
\end{defn}

\begin{defn}[Teorema]
!!^a{formula}~$!A$ merupakan suatu \emph{teorema} jika terdapat
!!a{derivation} atas $!A$ dari himpunan kosong. Kita menulis
$\Proves !A$ jika $!A$ merupakan teorema dan $\Proves/ !A$ jika bukan.
\end{defn}

\begin{defn}[Konsistensi]
Suatu himpunan !!{formula}s~$\Gamma$ \emph{konsisten} jika dan hanya
jika $\Gamma\Proves/ \lfalse$; jika tidak, himpunan itu
\emph{tak konsisten}.
\end{defn}

\begin{prop}[Refleksivitas]
\ollabel{prop:reflexivity}
Jika $!A \in \Gamma$, maka $\Gamma \Proves !A$.
\end{prop}

\begin{proof}
  Secara langsung, !!{formula}~$!A$ sendiri merupakan !!a{derivation} atas~$!A$
  dari~$\Gamma$.
\end{proof}

\begin{prop}[Monotonisitas]
\ollabel{prop:monotonicity}
Jika $\Gamma \subseteq \Delta$ dan $\Gamma \Proves !A$, maka $\Delta
\Proves !A$.
\end{prop}

\begin{proof}
Setiap !!{derivation} atas $!A$ dari $\Gamma$ juga merupakan
!!a{derivation} atas $!A$ dari~$\Delta$.
\end{proof}

\begin{prop}[Transitivitas]
\ollabel{prop:transitivity}
Jika $\Gamma \Proves !A$ dan $\{!A\} \cup \Delta \Proves
!B$, maka $\Gamma \cup \Delta \Proves !B$.
\end{prop}

\begin{proof}
  Misalkan $\{!A\} \cup \Delta \Proves !B$. Maka terdapat
  !!a{derivation} $!B_1$, \dots, $!B_l = !B$ dari~$\{!A\} \cup
  \Delta$. Beberapa langkah dalam !!{derivation} tersebut akan benar
  karena suatu aturan yang mengacu pada baris sebelumnya~$!B_i = !A$.
  Menurut hipotesis, terdapat !!a{derivation} atas~$!A$ dari~$\Gamma$,
  yakni !!a{derivation}~$!A_1$, \dots, $!A_k = !A$ dengan setiap
  $!A_i$ merupakan suatu aksioma, !!a{element} dari~$\Gamma$, atau
  benar menurut suatu aturan inferensi. Sekarang perhatikan urutan
  \[
  !A_1, \dots, !A_k = !A, !B_1, \dots, !B_l = !B.
  \]
  Urutan ini merupakan !!{derivation} yang benar atas~$!B$ dari
  $\Gamma \cup \Delta$, sebab setiap $!B_i = !A$ sekarang dibenarkan
  oleh pembenaran yang sama dengan yang membenarkan~$!A_k = !A$.
\end{proof}

Perhatikan bahwa, khususnya, ini berarti bahwa jika $\Gamma \Proves !A$
dan $!A \Proves !B$, maka $\Gamma \Proves !B$. Diperoleh pula bahwa
jika $!A_1, \dots, !A_n \Proves !B$ dan $\Gamma \Proves !A_i$ untuk
setiap~$i$, maka $\Gamma \Proves !B$.

\begin{prop}
\ollabel{prop:incons}
$\Gamma$ tak konsisten jika dan hanya jika $\Gamma \Proves !A$ untuk
setiap~$!A$.
\end{prop}

\begin{proof}
Latihan.
\end{proof}

\tagprob{FOL}
\begin{prob}
Buktikan \olref[fol][axd][ptn]{prop:incons}.
\end{prob}
\tagendprob

\tagprob{notFOL}
\begin{prob}
Buktikan \olref[pl][axd][ptn]{prop:incons}.
\end{prob}
\tagendprob

\begin{prop}[Kekompakan]
\ollabel{prop:proves-compact}
  \begin{enumerate}
  \item Jika $\Gamma \Proves !A$, maka terdapat himpunan bagian berhingga
    $\Gamma_0 \subseteq \Gamma$ sedemikian sehingga
    $\Gamma_0 \Proves !A$.
  \item Jika setiap himpunan bagian berhingga dari~$\Gamma$ konsisten,
    maka $\Gamma$ konsisten.
  \end{enumerate}
\end{prop}

\begin{proof}
  \begin{enumerate}
    \item Jika $\Gamma \Proves !A$, maka terdapat urutan berhingga
      !!{formula}s $!A_1$, \dots,~$!A_n$ sedemikian sehingga
      $!A \ident !A_n$ dan setiap $!A_i$ merupakan suatu aksioma logis,
      !!a{element} dari~$\Gamma$, atau mengikuti !!{formula}s sebelumnya
      melalui suatu aturan inferensi. Ambil $\Gamma_0$ sebagai himpunan semua
      $!A_i$ yang berada dalam~$\Gamma$. Maka !!{derivation} tersebut
      juga merupakan !!a{derivation} dari~$\Gamma_0$, sehingga
      $\Gamma_0 \Proves !A$.
    \item Ini adalah kontraposisi dari~(1) untuk kasus khusus
      $!A \ident \lfalse$.
\end{enumerate}
\end{proof}

\end{document}
